\section{「熊野寮なんかに入ったら就職活動にも響いて人生終わる」言説に対抗する}
\rightline{（文責：泉こなた（26卒））}

　こんにちは。26卒の泉こなたです。タイトルの通り、インターネット等でよく聞かれる「熊野寮なんかに入ったら就職活動にも響いて人生終わる」という言説に対抗しようと思います。

\subsection{はじめに}
　大前提として、私は基本的にはほとんどの進路選択を尊重しておりますので（スパイになる等は尊重しません）、4年で出て就職しようが、留年して就職しようが、研究者になろうが、定職に就かずダラダラとしていようが、革命家になろうがなんでもいいと思っています。そのうえで、上記の言説に対抗するために、いろいろと事実ベースで就職していく事例について書いていこうと思います。

\subsection{熊野寮生の就職事情}
　熊野寮生でもそれなりの人数が就活をし、就職先を得て退寮していきます。本節ではその就職先の中でも名の知れているものを具体的な企業名は伏せて紹介します。

～熊野寮生の就職先～
　・国家公務員（某省総合職）
　・某BIG4コンサル
　・某五大商社
　・某スーパーゼネコン
　・某自動車会社
　・某鉄道会社
　・某テレビ局
　・某証券会社
　・某メガバンク
　・某政府系金融機関
　・某保険会社

　ざっと挙げるだけでも、私が寮にいた数年でこれらに就職しています。一般的な京大生の就職先として見てみても、遜色ないように思います。

\subsection{なぜ熊野寮生は就職できるのか}
　上に挙げたような企業になぜ入ることができるのかを、ここでは考察していこうと思います。

\subsubsection{理由1：京大生はそもそも就活それなりに強い}
⇒熊野寮生も就職活動の場では「京大生」とみられることが多いです。京大生はそれなりに就職活動に強いので、熊野寮生も同様の状況に置かれます。某メガバンクなど、京大生をたくさん採用したいというような意識を持っている企業もそれなりにあります。

\subsubsection{理由2：寮生活がそのまま「ガクチカ」になる場合も}
⇒寮内組織の長をやったり、寮内の会議で議論の末議案の承認を得たりするという過程は、ほかの就活生がなかなかやっていないことです。そこらへんのサークルよりもすごいことをやっています。寮自治をガクチカに書く人は多くないので、エントリーシートの際に目に留まることも想定されます。寮自治に関わる中で自分が責任を持ってやり切ったことをきちんと伝えることができれば、まずよい評価を得られるでしょう。私も就職活動の際には自分のことを話す際の軸の1つに寮自治でいろいろとやってきたことを据えました。楽しく、責任を持ってやってきたことなので、すらすらと自分の言葉で話すこともできました。

\subsubsection{理由3：様々な人とコミュニケーションをとるので、面接での対応力が身につく}
⇒熊野寮には学部生から博士課程まで、年齢や性別、専門分野も様々住んでいます。多くの前提が異なるので、寮生と話す際には前提の違う人同士でも語ることのできるよう工夫する必要があります。また、熊野寮では様々な会議が行われるので、そこにおける議論や提起でも、参加していればスキルが磨かれます。就職活動で話すのは自分のことを知らない、前提の異なる人々です。そのような人々と相対しても、うまいことやっていける力が熊野寮では身につくのではないでしょうか。

\subsection{おわりに}
　本記事では熊野寮生の就活事情について述べました。私が言いたかったことをまとめると

\textbf{　1．熊野寮生は就職それなりにできます}

\textbf{　2．寮自治でいろいろとやってきたことが、就職活動に結果として生きるので、気にせず寮自治にゴリゴリ取り組もう}

　という具合です。最近は焦って1回生から就活するような人もちょいちょい見かけますが、まずは寮自治頑張ってみてもいいのかなと思います。こんなに裁量が与えられていろいろとできる環境はなかなかありません。退寮前に振り返ってみると、本当に良い場所だったなと思います。せっかく様々なインターネット言説を振り切って熊野寮に入るという決断をしてくれたならば、ぞんぶんに楽しみ尽くしてほしいところです。